\documentclass[11pt,a4paper]{article}
\usepackage[margin=2.2cm]{geometry}
\usepackage{ctex}
\usepackage{amsmath,amssymb,amsthm}
\usepackage{enumitem}
\usepackage{tasks}
\usepackage{array,tabularx}
\usepackage{multicol}
\usepackage{fancyhdr}
\usepackage{tikz}
\usepackage{graphicx}
\usepackage{xcolor}
\usepackage{needspace}

\pagestyle{fancy}
\fancyhf{}
\fancyhead[L]{\small 大模型训练与分布式计算 期末试卷}
\fancyhead[R]{\small 第 \thepage\ 页}
\renewcommand{\headrulewidth}{0.4pt}

\setlength{\parindent}{0pt}
\setlength{\parskip}{0pt}

% 工具命令
\newcommand{\blank}[1]{\underline{\hspace{#1}}}
\newcommand{\fn}[2]{\dfrac{#1}{#2}}

% 大题标题
\newcommand{\examsection}[2]{%
  \vspace{0.6cm}
  \noindent{\CJKfamily{SimHei}\bfseries\large #1\quad (#2)}\par
  \vspace{0.3cm}
}

% 题号 + 题干 一行写完（避免 parskip 陷阱）
\newcommand{\qitem}[2]{%
  \noindent\textbf{#1.}\;#2\par
  \vspace{0.3em}
}

% 选择题选项（4个一行/两行可控）
\newcommand{\fourchoices}[4]{%
  \begin{enumerate}[label=(\Alph*),leftmargin=2.2em,itemsep=0pt,parsep=0pt,topsep=0pt]
    \setlength{\itemindent}{0pt}
    \item #1
    \item #2
    \item #3
    \item #4
  \end{enumerate}
  \vspace{0.4em}
}

% 选择题选项（AB 一行，CD 一行；各占 1/2 宽度，区域内左对齐）
\newcommand{\fourchoicestwo}[4]{%
  \noindent
  \begin{minipage}[t]{0.5\linewidth}\makebox[2.2em][l]{(A)}#1\end{minipage}%
  \begin{minipage}[t]{0.5\linewidth}\makebox[2.2em][l]{(B)}#2\end{minipage}\par
  \noindent
  \begin{minipage}[t]{0.5\linewidth}\makebox[2.2em][l]{(C)}#3\end{minipage}%
  \begin{minipage}[t]{0.5\linewidth}\makebox[2.2em][l]{(D)}#4\end{minipage}\par
  \vspace{0.4em}
}

% 两栏式选项（短选项）
\newcommand{\twochoices}[4]{%
  \begin{tasks}[counter-format={(A)},label-align=left,label-offset={0.4em},label-width={1.6em},item-indent=2em,column-sep=1.5em](2)
    \task #1
    \task #2
    \task #3
    \task #4
  \end{tasks}
  \vspace{0.4em}
}

% 答题区横线
\newcommand{\answerlines}[1]{%
  \vspace{0.2cm}
  \foreach \i in {1,...,#1} {
    \noindent{\color{white}\rule{\linewidth}{0.4pt}}\par\vspace{0.5cm}
  }
}

\begin{document}

% ============ 试卷标题与说明 ============
\begin{center}
{\zihao{2}\CJKfamily{SimHei}\bfseries 大模型训练与分布式计算 期末试卷}\\[0.3cm]
{\zihao{4}\kaishu 满分 150 分 \quad 考试时间 150 分钟}
\end{center}

\vspace{1.5cm}
\noindent\rule{\linewidth}{1pt}

\vspace{0.2cm}
\noindent{\CJKfamily{SimHei}\bfseries 注意事项：}
\begin{enumerate}[leftmargin=2em,itemsep=0pt,parsep=0pt,topsep=0pt]
  \item 本试卷共 \textbf{三} 大题，包括选择题、填空题、解答题，共 \textbf{24} 小题。
  \item 请将答案写在答题区指定位置；解答题须写出必要的文字说明、推导过程或演算步骤。
  \item 涉及数学公式推导时，请清晰写出关键中间步骤；只写最终结果不得分。
  \item 本试卷中如无特殊说明，模型参数量记为 $\Phi$，隐藏维度记为 $h$，序列长度记为 $s$，批大小（Batch Size）记为 $b$，GPU 数量记为 $P$，混合精度训练默认使用 fp16/bf16 存储模型权重与梯度，优化器状态使用 fp32 存储。
\end{enumerate}

\vspace{0.2cm}
\noindent\rule{\linewidth}{0.6pt}

% ============ 一、选择题 ============
\examsection{一、选择题}{本大题共 12 小题，每小题 5 分，共 60 分。在每小题给出的四个选项中，只有一项是符合题目要求的。}

\qitem{1}{关于混合专家模型（Mixture of Experts, MoE）的基本概念，以下说法\textbf{正确}的是}{}
\fourchoices{MoE 中每个 Token 仅由一个专家处理，其余专家对该 Token 不产生任何影响}{在 Switch Transformer 中，使用 Top-1 路由意味着每个 Token 只被路由到一个专家}{MoE 的计算复杂度随专家数量线性增长，与激活 Token 数量无关}{MoE 通过增加专家数量可以无限扩大模型容量，而不会增加任何训练成本}

\qitem{2}{根据 Chinchilla Scaling Law（Hoffmann 等人提出的缩放定律），在给定计算量预算 $C$（单位：FLOPs）下，模型参数量 $N$ 与训练 Token 数 $D$ 的最优分配比例应满足 $C \approx 6ND$。当算力预算增加到原来的 $k$ 倍时，为了达到最优训练效果，模型参数量 $N$ 和训练数据量 $D$ 的理论增长比例大约为}{}
\fourchoices{$N$ 增长 $k^{0.73}$ 倍，$D$ 增长 $k^{0.27}$ 倍}{$N$ 增长 $k^{0.5}$ 倍，$D$ 增长 $k^{0.5}$ 倍}{$N$ 增长 $k^{0.3}$ 倍，$D$ 增长 $k^{0.7}$ 倍}{$N$ 增长 $k$ 倍，$D$ 保持不变}

\qitem{3}{在流水线并行（Pipeline Parallelism）中，采用 1F1B（One Forward, One Backward）调度策略。设流水线阶段数为 $P$，一个 Batch 内被切分成的微批（Micro-batch）数量为 $M$（满足 $M \ge P$）。若单次前向传播时间为 $F$，单次反向传播时间为 $B$，则在一个完整的训练步（Step）中，由于气泡（Bubble）导致的处理器空闲时间占总时间的比例（气泡率）为}{}
\fourchoicestwo{$\dfrac{P-1}{M+P-1}$}{$\dfrac{P-1}{M}$}{$\dfrac{M-1}{M+P-1}$}{$\dfrac{P}{M}$}
\newpage

\qitem{4}{在 Megatron-LM 样式的张量并行（Tensor Parallelism, TP）中，一个标准的 Transformer 层包含自注意力（Self-Attention）和多层感知机（MLP）两个子块。若 TP 的切分大小为 $t$，在单层前向传播过程中，各个 GPU 之间需要进行集体通信。如果不使用序列并行，该层前向传播总共需要执行的通信操作为}{}
\fourchoices{2 次 All-Reduce}{2 次 All-Gather 和 2 次 Reduce-Scatter}{4 次 All-Reduce}{1 次 All-Reduce 和 1 次 All-Gather}

\qitem{5}{在大模型训练中，激活重计算（Activation Checkpointing / Gradient Checkpointing）常用于以计算换显存。若对某一 Transformer 层采用选择性重计算（Selective Checkpointing）策略，丢弃其中 $75\%$ 容易重新计算的中间激活值。设该层原本的前向传播计算时间为 $T_f$，则采用该策略后，该层在反向传播时因重计算引入的额外时间开销约为}{}
\fourchoicestwo{$0.25 T_f$}{$0.5 T_f$}{$0.75 T_f$}{$1.5 T_f$}

\qitem{6}{关于 DeepSpeed 提出的 ZeRO（Zero Redundancy Optimizer）三种阶段（Stage 1/2/3）对模型状态显存的切分，下列说法\textbf{错误}的是}{}
\fourchoices{ZeRO-1 仅将优化器状态（Optimizer States）进行分片存储，不切分权重和梯度}{ZeRO-2 在 ZeRO-1 的基础上，进一步将梯度（Gradients）进行分片存储}{ZeRO-3 将优化器状态、梯度和模型权重（Weights）均进行分片存储}{ZeRO 各阶段通过分片存储激活值（Activations）来彻底解决超长文本训练时的 OOM 问题}

\qitem{7}{直接偏好优化（Direct Preference Optimization, DPO）是近年来替代传统基于 PPO 算法的 RLHF 的重要方法。已知提示词为 $x$，人类偏好的好回答为 $y_w$，较差回答为 $y_l$，参考模型为 $\pi_{ref}$，当前训练的策略模型为 $\pi_{\theta}$。DPO 损失函数（不含正则项）的形式为}{}
\fourchoices{$-\mathbb{E}_{(x, y_w, y_l)} \left[ \log \sigma \left( \beta \log \dfrac{\pi_\theta(y_w|x)}{\pi_{ref}(y_w|x)} - \beta \log \dfrac{\pi_\theta(y_l|x)}{\pi_{ref}(y_l|x)} \right) \right]$}{$-\mathbb{E}_{(x, y_w, y_l)} \left[ \log \sigma \left( \beta \log \dfrac{\pi_{ref}(y_w|x)}{\pi_\theta(y_w|x)} - \beta \log \dfrac{\pi_{ref}(y_l|x)}{\pi_\theta(y_l|x)} \right) \right]$}{$-\mathbb{E}_{(x, y_w, y_l)} \left[ \beta \log \dfrac{\pi_\theta(y_w|x)}{\pi_{ref}(y_l|x)} - \log \sigma \left( \dfrac{\pi_\theta(y_l|x)}{\pi_{ref}(y_w|x)} \right) \right]$}{$-\mathbb{E}_{(x, y_w, y_l)} \left[ \log \sigma \left( \log \dfrac{\pi_\theta(y_w|x)}{\pi_\theta(y_l|x)} - \beta \log \dfrac{\pi_{ref}(y_w|x)}{\pi_{ref}(y_l|x)} \right) \right]$}

\qitem{8}{在使用分布式数据并行（DDP）训练时，各节点需要对梯度进行同步。若集群有 $P$ 个 GPU 节点，梯度数据总大小为 $V$ 字节，采用环形全归约（Ring-AllReduce）算法。在每个训练步的梯度同步阶段，每个 GPU 发送的总数据量(单位：字节)为}{}
\fourchoicestwo{$V$}{$\dfrac{2(P-1)V}{P}$}{$\dfrac{(P-1)V}{P}$}{$2PV$}

\qitem{9}{在大模型预训练的初始阶段，学习率通常需要经历一个温热启动（Warm-up）阶段（即从 0 渐进增加到设定的最大学习率 $\eta_{max}$），其核心数学与物理机理是}{}
\fourchoices{随机初始化的权重导致早期梯度极不稳定，较小的学习率可以避免模型在训练初期发散}{增加前期梯度的绝对方差，使得模型更容易跳出局部最优解}{限制前期权重更新的 L2 范数，等效于引入强力的权重衰减（Weight Decay）}{缓解数据分布不均匀导致的早期过拟合}

\qitem{10}{对于一个标准的自回归 Decoder-only Transformer 模型（不使用激活重计算），在单步训练中，反向传播（Backward Pass）阶段所需的计算量（FLOPs）理论上大约是前向传播（Forward Pass）阶段的}{}
\fourchoicestwo{1 倍}{1.5 倍}{2 倍}{3 倍}

\qitem{11}{在人类反馈强化学习（RLHF）的 PPO 阶段，为了防止策略模型 $\pi_\theta$ 偏离初始的 SFT 模型 $\pi_{SFT}$ 太远导致"奖励作弊（Reward Hacking）"，通常会在奖励函数中引入 KL 散度惩罚。若原始奖励模型输出为 $R(x, y)$，惩罚系数为 $\beta > 0$，则经 KL 惩罚修正后的实际奖励值 $R_{modified}$ 表达式为}{}
\fourchoices{$R(x, y) - \beta \log \dfrac{\pi_{SFT}(y|x)}{\pi_\theta(y|x)}$}{$R(x, y) - \beta \log \dfrac{\pi_\theta(y|x)}{\pi_{SFT}(y|x)}$}{$R(x, y) + \beta \left( \pi_\theta(y|x) - \pi_{SFT}(y|x) \right)$}{$R(x, y) \cdot e^{-\beta D_{KL}(\pi_\theta \| \pi_{SFT})}$}

\qitem{12}{在 FP16 混合精度训练中，引入"动态损失缩放"（Dynamic Loss Scaling）机制是为了解决以下哪个数值计算瓶颈？}{}
\fourchoices{模型权重值过大导致溢出（Overflow）}{梯度数值太小，超出 FP16 的下限表示范围，导致数值下溢（Underflow）归零}{激活值在前向传播中由于矩阵乘法不断累加而造成的数值爆炸}{AdamW 优化器在一阶、二阶动量状态更新时的数值溢出}

% ============ 二、填空题 ============
\examsection{二、填空题}{本大题共 6 小题，每小题 5 分，共 30 分。把答案填在题中横线上。}

\qitem{13}{设某大模型训练任务的全局批大小（Global Batch Size）设为 1024。当前分配了 32 张 GPU 进行数据并行（DP）训练，若每张 GPU 因显存限制最大只能容纳微批大小（Micro-batch Size）为 4，则需要设置的梯度累积步数（Gradient Accumulation Steps）为 \blank{2cm}。}{}

\qitem{14}{某大模型拥有 $\Phi = 7 \times 10^9$（7B）参数。若使用 BF16 混合精度，并采用 AdamW 优化器（每个参数额外维护 FP32 版本的 Master Weight、一阶动量和二阶动量），在不考虑激活值、梯度和通信缓存的前提下，仅存储模型状态（Model States）本身所需的显存大小至少为 \blank{3cm} GB。（注：结果保留为整数，$1\text{ GB} = 10^9\text{ 字节}$）}{}

\qitem{15}{在 Megatron-LM 序列并行（Sequence Parallelism, SP）中，LayerNorm 和 Dropout 的激活值在序列长度维度 $S$ 上被切分。设张量并行大小为 $t$，Batch 大小为 $B$，隐藏维度为 $D$。在单层 Transformer 的前向传播中，由于引入序列并行，每个 GPU 在 SP 相关的通信中（All-Gather 和 Reduce-Scatter 各一次）需要通信的元素（Element）总数为 \blank{4cm}（用 $S, B, D, t$ 表示）。}{}

\qitem{16}{在 RLHF 偏好对齐中，若词表大小仅为 2（即只有 A, B 两个 Token），对于某特定 Prompt $x$，SFT 模型输出分布为 $\pi_{SFT}(A|x) = 0.5,\ \pi_{SFT}(B|x) = 0.5$；经过 PPO 训练后的策略模型分布为 $\pi_{\theta}(A|x) = 0.8,\ \pi_{\theta}(B|x) = 0.2$。则当前单步所产生的 KL 散度值 $D_{KL}(\pi_{\theta} \| \pi_{SFT})$ 为 \blank{2.5cm}。（取 $\ln 2 \approx 0.69,\ \ln 5 \approx 1.61$，结果保留两位小数）}{}

\qitem{17}{在因果语言模型（Causal Language Model）自回归预训练中，为了实现单次前向传播同时计算整个序列所有位置的 Loss，需要引入因果掩码（Causal Mask）。对于位置为 $i$ 的 Query 和位置为 $j$ 的 Key，当 $i < j$ 时，注意力得分矩阵（Attention Score）在 Softmax 之前需要被加上 \blank{3cm}。}{}

\qitem{18}{假设在一张单卡上训练某一模型时，其吞吐量为 $T$ tokens/s。现在使用 $N$ 张卡构建 DDP 集群进行分布式训练，由于节点间梯度同步通信导致了效率损耗，测得集群的弱扩展效率（Weak Scaling Efficiency）为 $90\%$。则该集群的整体实际训练吞吐量公式为 \blank{4cm}。}{}

% ============ 三、解答题 ============
\examsection{三、解答题}{本大题共 6 小题，共 60 分。解答应写出文字说明、证明过程或演算步骤。}

\qitem{19}{（本小题满分 10 分）}{}
\noindent 在大模型预训练中，"训练算力估算"是系统规划的核心。设一个 Decoder-only 架构的 Transformer 模型，其非嵌入层参数量为 $N$，隐藏维度为 $h$，前向传播中每个 Token 平均经历的操作可以近似简化为矩阵乘法（GEMM）。
\begin{enumerate}[label=(\arabic*),leftmargin=2em,itemsep=0.2em,parsep=0pt,topsep=0pt]
  \item 试从矩阵乘法的维度和乘加操作（MAC）的角度，推导证明：对单个 Token 进行前向传播所需的计算量大约为 $2N$ FLOPs；（4 分）
  \item 说明在训练阶段，为什么单步（一次 Forward 和一次 Backward）所需的总计算量大约为 $6N$ FLOPs/token；（4 分）
  \item 现有一款拥有 70B（$70 \times 10^9$）参数的模型，准备在含有 2T（$2 \times 10^{12}$）Tokens 的数据集上进行预训练。若使用一个 A100 GPU 集群进行训练，假设在混合精度和算子优化下，每张卡实际达到的训练吞吐算力（MFU）平均为 150 TFLOPs/s（$150 \times 10^{12}$ FLOPs/s），估算该预训练任务在 256 张 A100 上需要运行多少天？（2 分）
\end{enumerate}

\answerlines{10}

\needspace{6cm}
\qitem{20}{（本小题满分 10 分）}{}
\noindent DPO（Direct Preference Optimization）算法通过巧妙的数学变换，直接在偏好数据上优化语言模型，避免了 RLHF 中复杂的强化学习循环。在标准 RLHF 的对齐目标中，我们希望求解以下受约束的优化问题：
\[
\max_{\pi} \mathbb{E}_{x, y \sim \pi} [R(x, y)] - \beta D_{KL}(\pi(y|x) \parallel \pi_{ref}(y|x))
\]
其中 $\pi_{ref}$ 是 SFT 模型，$\pi$ 为待优化的策略模型，$R(x, y)$ 是学到的奖励函数，$\beta$ 是控制偏离度的常数。
\begin{enumerate}[label=(\arabic*),leftmargin=2em,itemsep=0.2em,parsep=0pt,topsep=0pt]
  \item 试证明：上述受约束优化问题的闭式最优解（Optimal Policy）具有如下形式：
  \[
  \pi^*(y|x) = \frac{1}{Z(x)} \pi_{ref}(y|x) \exp\left( \frac{R(x, y)}{\beta} \right)
  \]
  其中 $Z(x) = \sum_{y} \pi_{ref}(y|x) \exp\left( \frac{R(x, y)}{\beta} \right)$ 是归一化配分函数；（5 分）
  \item 结合 Bradley-Terry 偏好模型，即人类偏好 $y_w$ 优于 $y_l$ 的概率满足：$P(y_w \succ y_l \mid x) = \sigma(R(x, y_w) - R(x, y_l))$，利用（1）中的最优解关系，推导消去隐式奖励函数 $R(x, y)$，从而给出 DPO 损失函数的数学形式。（5 分）
\end{enumerate}

\answerlines{12}

\needspace{6cm}
\qitem{21}{（本小题满分 10 分）}{}
\noindent ZeRO（Zero Redundancy Optimizer）通过消除传统数据并行（DP）中的冗余状态，极大降低了单卡显存开销。设模型参数量为 $\Phi$，分布式数据并行大小（即 GPU 数量）为 $N_d$。使用 BF16 混合精度和 AdamW 优化器，不考虑激活值。
\begin{enumerate}[label=(\arabic*),leftmargin=2em,itemsep=0.2em,parsep=0pt,topsep=0pt]
  \item 写出在标准 DP（未应用 ZeRO）、ZeRO-1、ZeRO-2、ZeRO-3 下，每个 GPU 存储模型状态（参数、梯度、优化器状态）所需的显存大小（用 $\Phi$ 和 $N_d$ 表示的函数，单位：字节）；（4 分）
  \item 分析通信开销：标准 DP 每一训练步通过 All-Reduce 进行梯度同步，通信量（每个 GPU 移入/移出数据量）为 $2\Phi$ 参数个元素。请详细说明 ZeRO-3 在前向传播（Forward）和反向传播（Backward）中是如何通过集体通信更新和释放权重的，并证明其单步总通信量为标准 DP 的 1.5 倍。（6 分）
\end{enumerate}

\answerlines{5}

\needspace{6cm}
\qitem{22}{（本小题满分 10 分）}{}
\noindent 在优化器设计中，AdamW（Decoupled Weight Decay）相比传统的 Adam + L2 正则化在训练大模型时具有更好的泛化性能。
\begin{enumerate}[label=(\arabic*),leftmargin=2em,itemsep=0.2em,parsep=0pt,topsep=0pt]
  \item 设当前步梯度为 $g_t$，一阶动量为 $m_t$，二阶动量为 $v_t$，学习率为 $\eta_t$，权重衰减系数为 $\lambda$。分别写出 Adam 结合 L2 正则化（即把 $\frac{\lambda}{2} \|\theta\|_2^2$ 显式加进 Loss 中）以及 AdamW 在第 $t$ 步的参数更新方程；（4 分）
  \item 结合上述更新方程，从数学上推导并阐述：为什么在 Adam 优化器中，传统的 L2 正则化机制不能等效于 SGD 中的权重衰减（Weight Decay）？而 AdamW 是如何纠正这一偏差的？（6 分）
\end{enumerate}

\answerlines{5}

\needspace{6cm}
\qitem{23}{（本小题满分 10 分）}{}
\noindent 在大规模大模型训练中，通常需要将张量并行（TP）、流水线并行（PP）和数据并行（DP）结合，形成"3D 混合并行"。现有一个包含 $G = 32$ 张 GPU 的高带宽互联集群，计划进行 3D 并行配置。
\begin{enumerate}[label=(\arabic*),leftmargin=2em,itemsep=0.2em,parsep=0pt,topsep=0pt]
  \item 给出 3D 并行的配置约束等式。若设置张量并行大小为 $t = 4$，流水线并行大小为 $p = 4$，则数据并行大小 $d$ 应为多少？（2 分）
  \item 为了兼顾通信效率，GPU 分配应当考虑物理拓扑（如单机内 8 卡通过 NVLink 高速互联，机间通过低带宽 InfiniBand 互联）。若有 4 台单机 8 卡物理节点（共 32 卡），GPU 物理编号为 0 到 31。请给出一套合理的 GPU 并行组划分方案，使 TP 组内通信全部限制在物理单机内部，并列出编号为 0 的 GPU 所在的 TP 组、PP 组和 DP 组包含的所有 GPU 编号；（4 分）
  \item 简述 3D 并行中，为何通常将张量并行限制在单机内部，而流水线并行和数据并行可以跨机布置？（4 分）
\end{enumerate}

\answerlines{3}

\needspace{6cm}
\qitem{24}{（本小题满分 10 分）}{}
\noindent 长文本（Long Context）训练是当前大模型的核心竞争点之一。环形注意力（Ring Attention）通过在环形拓扑的 GPU 链上切分并循环传递 Key-Value 块，成功突破了单卡处理超长序列的内存墙。设待训练的序列长度为 $S$，隐藏维度为 $D$，注意力头数为 $H$。我们使用由 $N_{ring}$ 张 GPU 组成的环形通信拓扑。
\begin{enumerate}[label=(\arabic*),leftmargin=2em,itemsep=0.2em,parsep=0pt,topsep=0pt]
  \item 详细描述 Ring Attention 的并行工作原理和数据流向：如何通过非阻塞对等通信（Non-blocking P2P Send/Recv），将局部计算与 KV 块的传输进行重叠（Overlap），从而实现无全局 A-matrix（即无 $S \times S$ 完整注意力矩阵）显存驻留的因果注意力计算；（4 分）
  \item 设 $S = 1{,}048{,}576$（1M），$D = 8192$，$H = 64$。若不采用 FlashAttention 也不做任何分块，直接在单张卡上计算并存储单层的 Softmax 注意力分数矩阵，假设每个元素使用 FP32（4 字节）存储，计算单层单头注意力矩阵所需的理论显存大小。这说明了什么问题？（3 分）
  \item 结合 FlashAttention（在 SRAM 中完成局部块的 Softmax 和加权累加），推导说明在使用 Ring Attention 后，每张 GPU 在计算注意力时，其显存中与注意力计算相关的核心激活值显存占用（Q、K、V 激活值）的复杂度成功降低到了 $O\left( \frac{S}{N_{ring}} \cdot D \right)$。若 $N_{ring} = 32$，分析该配置下每张卡的最大文本处理上限能达到怎样的规模？（3 分）
\end{enumerate}

\answerlines{12}

% ============ 答案另起一页 ============
\newpage
\begin{center}
{\zihao{2}\heiti 参考答案与解析}\\[0.3cm]
{\small 命题：大模型训练与分布式计算课程组 \quad 校对：教研组}
\end{center}
\vspace{0.3cm}
\noindent\rule{\linewidth}{0.6pt}

\subsection*{\heiti 一、选择题答案}

\begin{multicols}{4}
\begin{enumerate}[label=\textbf{\arabic*.},leftmargin=*,itemsep=0.1em,parsep=0pt,topsep=0pt]
  \item B
  \item B
  \item A
  \item A
  \item C
  \item D
  \item A
  \item B
  \item A
  \item C
  \item B
  \item B
\end{enumerate}
\end{multicols}

\vspace{0.3cm}
\noindent\textbf{1. B}\quad (A) 错误：标准 MoE 常采用 Top-$k$ 路由（$k>1$），一个 Token 可同时被多个专家处理；即便在 Top-1 情形下，未被选中的专家虽不直接参与该 Token 的前向计算，路由网络与负载均衡机制仍会影响整体训练。(B) 正确：Switch Transformer 明确采用 Top-1 路由，每个 Token 经路由器打分后仅分配给得分最高的一个专家。(C) 错误：稀疏 MoE 的单步计算量主要取决于被激活的专家数（通常为 Top-$k$）与 Token 数量，而非专家总数 $E$。(D) 错误：增加专家会扩大总参数量及存储、通信与负载均衡开销，并非不增加任何训练成本。故选 B。

\noindent\textbf{2. B}\quad Chinchilla Scaling Law 证明，在算力增加时，参数量 $N$ 和数据集 Token 数 $D$ 应当以大致相等的比例缩放（即 $N\propto C^{0.5}$，$D\propto C^{0.5}$）。旧的 Kaplan 定律认为参数量增长应快于数据量（$N\propto C^{0.73}$），Chinchilla 修正了这一结论。故选 B。

\noindent\textbf{3. A}\quad 在 1F1B 流水线并行中，前向气泡和反向气泡集中在首尾。一个 Step 中流水线气泡部分等效于 $P-1$ 个微批的前向和反向传播时间，非气泡有效计算时间为 $M(F+B)$。总运行时间为 $(M+P-1)(F+B)$，气泡率为 $\dfrac{P-1}{M+P-1}$。故选 A。

\noindent\textbf{4. A}\quad 标准 Megatron 张量并行中：自注意力块在 $O$ 矩阵行切分时需要 1 次 All-Reduce；MLP 块第二层行切分结束时需要 1 次 All-Reduce。单层前向传播共需 2 次 All-Reduce。故选 A。

\noindent\textbf{5. C}\quad 丢弃 75\% 的激活值意味着反向传播时需对这 75\% 部分重新进行一次前向计算以恢复激活值，额外引入的时间开销约为 $0.75 T_f$。故选 C。

\noindent\textbf{6. D}\quad ZeRO 解决的是模型状态（Model States）的冗余存储问题，不对激活值做物理维度的切分。长文本导致的 OOM 主要由激活值过大引起，通常依靠序列并行（SP）、FlashAttention 或环形注意力（Ring Attention）解决。故选 D。

\noindent\textbf{7. A}\quad DPO 通过对 RLHF 的闭式解求反函数，推导出的损失函数为：
\[
\mathcal{L}_{DPO}(\theta;\pi_{ref})=-\mathbb{E}_{(x,y_w,y_l)}\!\left[\log\sigma\!\left(\beta\log\frac{\pi_\theta(y_w|x)}{\pi_{ref}(y_w|x)}-\beta\log\frac{\pi_\theta(y_l|x)}{\pi_{ref}(y_l|x)}\right)\right]
\]
该损失使模型提高偏好样本 $y_w$ 概率的同时降低非偏好样本 $y_l$ 的概率，并以 $\pi_{ref}$ 作为隐式 KL 散度约束。故选 A。

\noindent\textbf{8. B}\quad Ring-AllReduce 包含 Scatter-Reduce 与 All-Gather 两个阶段，每阶段每个节点传送 $P-1$ 次大小为 $V/P$ 的数据包。每个 GPU 总发送数据量为 $2\times(P-1)\times\dfrac{V}{P}=\dfrac{2(P-1)V}{P}$ 字节。故选 B。

\noindent\textbf{9. A}\quad 训练初期模型权重随机初始化，输出分布混乱，交叉熵损失梯度可能非常大且方向多变。若一开始就使用较大的 $\eta_{\max}$，容易在最初几个 Step 摧毁合理的权重初始化，导致梯度爆炸或数值溢出而训练发散。Warm-up 通过较小步长让模型逐渐适应数据分布，使梯度稳定。故选 A。

\noindent\textbf{10. C}\quad 前向传播执行一次矩阵乘法 $Y=XW$（$2MND$ FLOPs）。反向传播需执行两次同等规模的矩阵乘法：关于输入激活值的梯度 $\dfrac{\partial L}{\partial X}=\dfrac{\partial L}{\partial Y}W^{\top}$，以及关于权重的梯度 $\dfrac{\partial L}{\partial W}=X^{\top}\dfrac{\partial L}{\partial Y}$。因此反向是前向的 2 倍，总训练步为 3 倍。故选 C。

\noindent\textbf{11. B}\quad KL 散度 $D_{KL}(\pi_\theta\parallel\pi_{SFT})=\mathbb{E}_{y\sim\pi_\theta}\!\left[\log\dfrac{\pi_\theta(y|x)}{\pi_{SFT}(y|x)}\right]$。在 PPO 目标中希望最大化修正后的奖励并减去 KL 惩罚：$\text{Objective}=R(x,y)-\beta\log\dfrac{\pi_\theta(y|x)}{\pi_{SFT}(y|x)}$，超出 $\pi_{SFT}$ 分布越多的 token 惩罚越大，从而保持输出自然度。故选 B。

\noindent\textbf{12. B}\quad FP16 数值表示范围很窄（最小正非零正常数约 $6.10\times10^{-5}$）。训练深层网络时许多梯度绝对值常落在 $[10^{-8},10^{-5}]$，直接使用 FP16 会因低于精度下限被截断为 0（数值下溢）。动态损失缩放在前向计算完 Loss 后乘以大因子 $S$（如 65536），使反向梯度放大到 FP16 可表示范围，最后更新权重前再除以 $S$。故选 B。

\subsection*{\heiti 二、填空题答案}

\noindent\textbf{13.}\quad $8$\quad 全局批大小满足 $\text{Global Batch Size}=\text{GPU 数量}\times\text{微批大小}\times\text{梯度累积步数}$，代入 $1024=32\times4\times\text{Steps}$，得 $\text{Steps}=8$。

\noindent\textbf{14.}\quad $112$\quad BF16 混合精度下，每个参数占用模型状态显存：权重 $2$ 字节 + 梯度 $2$ 字节 + AdamW 优化器状态（FP32 Master Weight、一阶动量、二阶动量各 $4$ 字节）$12$ 字节，合计 $16$ 字节。7B 模型最小模型状态显存 $M=7\times10^{9}\times16=112\times10^{9}$ 字节 $=112$ GB。

\noindent\textbf{15.}\quad $2\dfrac{t-1}{t}SBD$\quad LN/Dropout 后通过 All-Gather 拼接序列，传输 $\dfrac{t-1}{t}SBD$ 元素；自注意力输出后通过 Reduce-Scatter 切分回各 GPU，同样传输 $\dfrac{t-1}{t}SBD$ 元素，累计 $2\dfrac{t-1}{t}SBD$。

\noindent\textbf{16.}\quad $0.19$\quad $D_{KL}=0.8\ln\dfrac{0.8}{0.5}+0.2\ln\dfrac{0.2}{0.5}=0.8\ln(1.6)+0.2\ln(0.4)\approx0.8\times0.4700+0.2\times(-0.9163)=0.376-0.183\approx0.19$。

\noindent\textbf{17.}\quad $-\infty$（或 $-10^9$ 等极小负数）\quad 当 $i<j$ 时，Query 处于当前 Token 之前，自回归预测不能看到未来 Token，故在 Softmax 前将该注意力得分加 $-\infty$，使其 Softmax 后权重严格为 0。

\noindent\textbf{18.}\quad $0.9NT$\quad 弱扩展效率为实际总吞吐量与单卡无损理想吞吐量的比例，效率为 $90\%$ 且有 $N$ 张卡时，$\text{Throughput}=N\cdot T\cdot 90\%=0.9NT$。

\subsection*{\heiti 三、解答题答案}

\subsubsection*{19. 训练算力估算（10 分）}

\noindent\textbf{(1) 前向传播计算量约为 $2N$ FLOPs（4 分）}

单层 Transformer 中主要计算集中在各种投影矩阵上，设隐藏维度为 $h$：
\begin{itemize}[leftmargin=2em,itemsep=0pt,parsep=0pt,topsep=0pt]
  \item 注意力 $Q,K,V$ 投影矩阵大小均为 $[h,h]$，共 $3h^2$ 参数，对应 $3h^2$ MACs；输出投影 $O$ 矩阵 $[h,h]$，对应 $h^2$ MACs。
  \item MLP 模块门控版本包含两个升维投影（$h\to 4h$）和一个降维投影（$4h\to h$），参数量约 $2(h\times4h)+(4h\times h)=12h^2$（非门控版本为 $8h^2$），对应 $12h^2$ MACs。
  \item 1 次 MAC = 1 次乘法 + 1 次加法 = 2 FLOPs。
\end{itemize}
汇总单层非嵌入参数量 $N_{\text{layer}}\approx 12h^2$，对应 $12h^2$ MACs $=24h^2$ FLOPs $=2N_{\text{layer}}$ FLOPs。推广至 $L$ 层全模型，单 Token 前向计算量约为 $2N$ FLOPs。（证明毕）

\noindent\textbf{(2) 训练单步约为 $6N$ FLOPs/token（4 分）}

\begin{itemize}[leftmargin=2em,itemsep=0pt,parsep=0pt,topsep=0pt]
  \item 前向传播：$2N$ FLOPs/token。
  \item 反向传播需计算两部分梯度：关于上一层激活值的梯度（与前向同规模，$2N$ FLOPs/token），关于该层参数权重的梯度（与前向同规模，$2N$ FLOPs/token），合计 $4N$ FLOPs/token。
\end{itemize}
单步总计算量 $=$ Forward $+$ Backward $=2N+4N=6N$ FLOPs/token。（证明毕）

\noindent\textbf{(3) 训练天数估算（2 分）}

总训练计算量：
\[
C_{\text{total}}=6\times N\times D=6\times(70\times10^{9})\times(2\times10^{12})=8.4\times10^{23}\text{ FLOPs}.
\]
单张 A100 每日吞吐：$150\times10^{12}\times86400=1.296\times10^{19}$ FLOPs/day。256 张 A100 每日总吞吐：$256\times1.296\times10^{19}\approx3.318\times10^{21}$ FLOPs/day。
\[
\text{Days}=\frac{8.4\times10^{23}}{3.318\times10^{21}}\approx253.2\text{ 天}.
\]
在 256 张 A100 上大约需要运行 253 天。

\subsubsection*{20. DPO 数学推导（10 分）}

\noindent\textbf{(1) 闭式最优解（5 分）}

优化目标为 $\max_{\pi}\sum_{y}\pi(y|x)R(x,y)-\beta\sum_{y}\pi(y|x)\log\dfrac{\pi(y|x)}{\pi_{ref}(y|x)}$，约束 $\sum_{y}\pi(y|x)=1$。重写为：
\[
\max_{\pi}\beta\sum_{y}\pi(y|x)\!\left[\log\!\left(\exp\!\left(\frac{R(x,y)}{\beta}\right)\pi_{ref}(y|x)\right)-\log\pi(y|x)\right].
\]
定义合法概率分布 $\pi^*(y|x)=\dfrac{1}{Z(x)}\pi_{ref}(y|x)\exp\!\left(\dfrac{R(x,y)}{\beta}\right)$，其中 $Z(x)=\sum_{y}\pi_{ref}(y|x)\exp\!\left(\dfrac{R(x,y)}{\beta}\right)$，则 $\pi_{ref}(y|x)\exp\!\left(\dfrac{R(x,y)}{\beta}\right)=Z(x)\pi^*(y|x)$，代回：
\[
\max_{\pi}\beta\log Z(x)\underbrace{\sum_{y}\pi(y|x)}_{=1}-\beta\sum_{y}\pi(y|x)\log\frac{\pi(y|x)}{\pi^*(y|x)}
=\beta\log Z(x)-\beta D_{KL}(\pi(y|x)\parallel\pi^*(y|x)).
\]
由于 $\beta\log Z(x)$ 与 $\pi$ 无关且 $D_{KL}(\pi\parallel\pi^*)\ge 0$，最大化当且仅当 KL 散度为 0，即 $\pi(y|x)=\pi^*(y|x)=\dfrac{1}{Z(x)}\pi_{ref}(y|x)\exp\!\left(\dfrac{R(x,y)}{\beta}\right)$。（证明毕）

\noindent\textbf{(2) DPO 损失推导（5 分）}

对最优解取对数：$\log\dfrac{\pi^*(y|x)}{\pi_{ref}(y|x)}=\dfrac{R(x,y)}{\beta}-\log Z(x)$，得 $R(x,y)=\beta\log\dfrac{\pi^*(y|x)}{\pi_{ref}(y|x)}+\beta\log Z(x)$。代入 Bradley-Terry 模型作差：
\[
R(x,y_w)-R(x,y_l)=\beta\log\frac{\pi^*(y_w|x)}{\pi_{ref}(y_w|x)}-\beta\log\frac{\pi^*(y_l|x)}{\pi_{ref}(y_l|x)},
\]
配分函数项 $\beta\log Z(x)$ 被完全消去。用当前策略 $\pi_\theta$ 替换 $\pi^*$：
\[
P(y_w\succ y_l\mid x)=\sigma\!\left(\beta\log\frac{\pi_\theta(y_w|x)}{\pi_{ref}(y_w|x)}-\beta\log\frac{\pi_\theta(y_l|x)}{\pi_{ref}(y_l|x)}\right).
\]
对该偏好数据集极大似然估计，负对数似然损失即 DPO 损失：
\[
\mathcal{L}_{DPO}(\theta)=-\mathbb{E}_{(x,y_w,y_l)}\!\left[\log\sigma\!\left(\beta\log\frac{\pi_\theta(y_w|x)}{\pi_{ref}(y_w|x)}-\beta\log\frac{\pi_\theta(y_l|x)}{\pi_{ref}(y_l|x)}\right)\right].
\]
（推导毕）

\subsubsection*{21. ZeRO 显存与通信分析（10 分）}

\noindent\textbf{(1) 各阶段单卡显存（4 分）}

混合精度下权重 $2\Phi$ 字节、梯度 $2\Phi$ 字节、优化器状态 $12\Phi$ 字节。
\begin{itemize}[leftmargin=2em,itemsep=0pt,parsep=0pt,topsep=0pt]
  \item 标准 DP：$M_{DP}=2\Phi+2\Phi+12\Phi=16\Phi$ 字节。
  \item ZeRO-1（仅分片优化器状态）：$M_{\text{ZeRO1}}=4\Phi+\dfrac{12\Phi}{N_d}$ 字节。
  \item ZeRO-2（分片梯度与优化器状态）：$M_{\text{ZeRO2}}=2\Phi+\dfrac{14\Phi}{N_d}$ 字节。
  \item ZeRO-3（分片全部模型状态）：$M_{\text{ZeRO3}}=\dfrac{16\Phi}{N_d}$ 字节。
\end{itemize}

\noindent\textbf{(2) ZeRO-3 通信量为标准 DP 的 1.5 倍（6 分）}

ZeRO-3 通信工作流程：
\begin{itemize}[leftmargin=2em,itemsep=0pt,parsep=0pt,topsep=0pt]
  \item 前向传播：每 GPU 只保存 $\dfrac{1}{N_d}$ 权重，计算到某层时发起 All-Gather 收集完整权重，计算后立即释放。整网前向共进行一次全模型权重 All-Gather，通信量 $\Phi$。
  \item 反向传播：再次 All-Gather 收集该层权重（通信量 $\Phi$），计算梯度后释放；算出的完整梯度通过 Reduce-Scatter 归约并分片存储（通信量 $\Phi$），随后释放非本分片梯度。
\end{itemize}
通信量对比：
\begin{itemize}[leftmargin=2em,itemsep=0pt,parsep=0pt,topsep=0pt]
  \item 标准 DP：反向结束时一次 All-Reduce，等价于 Reduce-Scatter + All-Gather，总通信量 $C_{DP}=2\Phi$。
  \item ZeRO-3：前向 All-Gather（$\Phi$）+ 反向 All-Gather（$\Phi$）+ 反向 Reduce-Scatter（$\Phi$），总通信量 $C_{\text{ZeRO3}}=3\Phi$。
  \item 比例：$\dfrac{C_{\text{ZeRO3}}}{C_{DP}}=\dfrac{3\Phi}{2\Phi}=1.5$。
\end{itemize}
即 ZeRO-3 单步总通信量为标准 DP 的 1.5 倍。（证明毕）

\subsubsection*{22. AdamW 与 L2 正则化（10 分）}

\noindent\textbf{(1) 更新方程（4 分）}

\textbf{Adam 结合 L2 正则化}：损失 $\tilde{\mathcal{L}}(\theta)=\mathcal{L}(\theta)+\dfrac{\lambda}{2}\|\theta\|_2^2$，梯度 $\tilde{g}_t=\nabla_\theta\mathcal{L}(\theta_t)+\lambda\theta_t$。
\[
m_t=\beta_1 m_{t-1}+(1-\beta_1)\tilde{g}_t,\quad v_t=\beta_2 v_{t-1}+(1-\beta_2)\tilde{g}_t^2,
\]
\[
\hat{m}_t=\frac{m_t}{1-\beta_1^t},\quad\hat{v}_t=\frac{v_t}{1-\beta_2^t},\quad\theta_{t+1}=\theta_t-\frac{\eta_t}{\sqrt{\hat{v}_t}+\epsilon}\hat{m}_t.
\]
\textbf{AdamW}：梯度 $g_t=\nabla_\theta\mathcal{L}(\theta_t)$（不含正则项），一、二阶动量用 $g_t$ 更新（公式同上），最后将权重衰减独立出来：
\[
\theta_{t+1}=\theta_t-\eta_t\lambda\theta_t-\frac{\eta_t}{\sqrt{\hat{v}_t}+\epsilon}\hat{m}_t.
\]

\noindent\textbf{(2) 推导与原理解释（6 分）}

在 SGD 中引入 L2 正则：$\theta_{t+1}=\theta_t-\eta_t(g_t+\lambda\theta_t)=(1-\eta_t\lambda)\theta_t-\eta_t g_t$，与直接权重衰减数学上完全等价。

但在 Adam 中，梯度更新需除以二阶动量历史 $\sqrt{\hat{v}_t}$。若使用 L2 正则，代入 $\hat{m}_t$ 中的正则项同样被二阶动量缩放：
\[
\theta_{t+1}\approx\theta_t-\frac{\eta_t}{\sqrt{\hat{v}_t}+\epsilon}\big(\nabla\mathcal{L}(\theta_t)+\lambda\theta_t\big).
\]
这意味着：对历史梯度大、$\sqrt{\hat{v}_t}$ 大的权重，实际权重衰减系数变为 $\dfrac{\lambda}{\sqrt{\hat{v}_t}}$，衰减强度被严重削弱；对历史梯度极小、$\sqrt{\hat{v}_t}$ 极小的权重，衰减强度反而被放大。导致权重衰减不均，正则化效果大打折扣。

AdamW 的纠正：将 $\lambda\theta_t$ 从动量更新中剥离，直接置于最后更新等式中，保证无论参数历史梯度如何，每 step 的权重衰减比例都严格恒定为 $\eta_t\lambda$，不受二阶动量干扰，从而完美恢复正则化初衷，在大模型训练中表现出更佳的泛化性和数值稳定性。

\subsubsection*{23. 3D 混合并行（10 分）}

\noindent\textbf{(1) 配置约束与 $d$ 的取值（2 分）}

3D 并行约束等式 $G=t\times p\times d$。代入 $G=32$，$t=4$，$p=4$：$32=4\times4\times d\Rightarrow 16d=32\Rightarrow d=2$。数据并行大小 $d$ 应为 2。

\noindent\textbf{(2) GPU 并行组划分（4 分）}

为使 TP 通信全部局限在单机内部，TP 组必须占用同一物理机器内的连续卡。4 台机器 Node 0--3 各含 8 卡，编号 0--7、8--15、16--23、24--31。因 $t=4$，每台机器容纳 $8/4=2$ 个 TP 组：
\begin{itemize}[leftmargin=2em,itemsep=0pt,parsep=0pt,topsep=0pt]
  \item TP 组 0：[0,1,2,3]（Node 0）；TP 组 1：[4,5,6,7]（Node 0）；依此类推。
\end{itemize}
GPU 0 所在并行组：
\begin{itemize}[leftmargin=2em,itemsep=0pt,parsep=0pt,topsep=0pt]
  \item TP 组：[0,1,2,3]。
  \item PP 组（跨阶段连接）：[0,4,8,12]（跨机布局保证阶段间对应拓扑）。
  \item DP 组（同阶段处理不同数据）：[0,16]。
\end{itemize}

\noindent\textbf{(3) TP 限单机、PP/DP 跨机的原因（4 分）}

\begin{itemize}[leftmargin=2em,itemsep=0pt,parsep=0pt,topsep=0pt]
  \item 张量并行（TP）：切分算子内部矩阵，每次前向/反向都需高频 All-Reduce（每个 Transformer 层数次），对延迟和带宽要求极高，跨机网络瓶颈会立刻导致 GPU 严重空闲，故须限制在单机内（利用机内 NVLink 吞吐）。
  \item 流水线并行（PP）：通信只发生在 Stage 边界（传输层间边界激活值），一个 Batch 期间通信频率远低于 TP，每次数据量较小，跨机通信对整体性能影响较低。
  \item 数据并行（DP）：通信只在反向传播结束后梯度同步阶段进行（通过 Bucket 化 All-Reduce 异步重叠），随 Batch Size 变大通信频率极低，可通过 IB 跨机网络轻松承载。
\end{itemize}

\subsubsection*{24. Ring Attention 长文本训练（10 分）}

\noindent\textbf{(1) 工作原理与数据流（4 分）}

\begin{itemize}[leftmargin=2em,itemsep=0pt,parsep=0pt,topsep=0pt]
  \item \textbf{空间切分}：将长度 $S$ 的输入序列在序列维度切分为 $N_{ring}$ 等份，每张 GPU 仅保留 $S_{\text{chunk}}=\dfrac{S}{N_{ring}}$ 的 $Q,K,V$ 块。
  \item \textbf{环形通信与局部计算重叠}：GPU $i$ 先计算本地 $Q_i$ 与本地 $K_i$ 的注意力并累加 $V_i$；同时启动非阻塞 P2P 发送，将当前 $K_i,V_i$ 发给环中下一个 GPU $i+1$，并从前一个 GPU $i-1$ 异步接收 $K,V$；传输的同时 GPU $i$ 用本地 $Q_i$ 与已接收的旧 $K,V$ 继续注意力计算。该过程循环 $N_{ring}-1$ 次，直至所有 $Q$ 与全网 $K,V$ 完成计算。
  \item \textbf{无全局矩阵驻留}：分块层面计算，GPU 只需在本地 SRAM 中计算局部 Softmax 和局部加权和，从不在任何 GPU 的 HBM 中实例化完整 $S\times S$ 注意力矩阵，规避长文本训练显存爆炸问题。
\end{itemize}

\noindent\textbf{(2) 单卡显存开销（3 分）}

单层单头 Softmax 注意力矩阵大小 $S\times S=1{,}048{,}576\times1{,}048{,}576$ 元素，元素总数 $S^2\approx1.10\times10^{12}$。FP32 每元素 4 字节：
\[
\text{Memory}=1.10\times10^{12}\times4\approx4.4\times10^{12}\text{ 字节}\approx4400\text{ GB}=4.4\text{ TB}.
\]
单卡不分块直接计算存储 1M 长度注意力矩阵需高达 4.4 TB 显存，远超当前单张 A100（80GB）容量，说明不进行分块和并行切分的长文本注意力计算在物理上完全无法实现。

\noindent\textbf{(3) 复杂度与处理上限（3 分）}

在 FlashAttention 与 Ring Attention 结合下，每张 GPU 仅持有长度 $\dfrac{S}{N_{ring}}$ 的序列片段。HBM 常驻激活值仅包含：本地 $Q$ 块（$\dfrac{S}{N_{ring}}\times D$）、当前计算与接收中的 $K,V$ 双缓冲（$2\times\dfrac{S}{N_{ring}}\times D$）、最终输出 $O$（$\dfrac{S}{N_{ring}}\times D$）。核心激活值显存复杂度被压缩到 $O\!\left(\dfrac{S}{N_{ring}}\cdot D\right)$ 的线性级别。

若 $N_{ring}=32$，处理 $S=1\text{M}$ 文本时每张 GPU 等效本地长度 $S_{\text{chunk}}=\dfrac{10^6}{32}\approx31250$。考虑 $D=8192$，这部分常驻激活值显存仅在数 GB 级别（相比 4.4 TB 降低几个数量级）。在 A100-80GB 留有 40GB 激活值预算下，得益于 $O\!\left(\dfrac{S}{N_{ring}}\cdot D\right)$ 的线性扩展性，32 卡集群理论极限处理长度可轻松突破 2M 至 4M 的极长文本规模。

\end{document}